% Options for packages loaded elsewhere
\PassOptionsToPackage{unicode}{hyperref}
\PassOptionsToPackage{hyphens}{url}
\documentclass[
  11pt,
]{article}
\usepackage{xcolor}
\usepackage[margin=25mm]{geometry}
\usepackage{amsmath,amssymb}
\setcounter{secnumdepth}{-\maxdimen} % remove section numbering
\usepackage{iftex}
\ifPDFTeX
  \usepackage[T1]{fontenc}
  \usepackage[utf8]{inputenc}
  \usepackage{textcomp} % provide euro and other symbols
\else % if luatex or xetex
  \usepackage{unicode-math} % this also loads fontspec
  \defaultfontfeatures{Scale=MatchLowercase}
  \defaultfontfeatures[\rmfamily]{Ligatures=TeX,Scale=1}
\fi
\usepackage{lmodern}
\ifPDFTeX\else
  % xetex/luatex font selection
\fi
% Use upquote if available, for straight quotes in verbatim environments
\IfFileExists{upquote.sty}{\usepackage{upquote}}{}
\IfFileExists{microtype.sty}{% use microtype if available
  \usepackage[]{microtype}
  \UseMicrotypeSet[protrusion]{basicmath} % disable protrusion for tt fonts
}{}
\makeatletter
\@ifundefined{KOMAClassName}{% if non-KOMA class
  \IfFileExists{parskip.sty}{%
    \usepackage{parskip}
  }{% else
    \setlength{\parindent}{0pt}
    \setlength{\parskip}{6pt plus 2pt minus 1pt}}
}{% if KOMA class
  \KOMAoptions{parskip=half}}
\makeatother
\setlength{\emergencystretch}{3em} % prevent overfull lines
\providecommand{\tightlist}{%
  \setlength{\itemsep}{0pt}\setlength{\parskip}{0pt}}
\usepackage{bookmark}
\IfFileExists{xurl.sty}{\usepackage{xurl}}{} % add URL line breaks if available
\urlstyle{same}
\hypersetup{
  hidelinks,
  pdfcreator={LaTeX via pandoc}}

\author{}
\date{}

\begin{document}

\section{Thom isomorphisms and K-orientations in
KK}\label{thom-isomorphisms-and-k-orientations-in-kk}

\emph{Written by GPT-6.1 Sol (OpenAI). Public domain (CC0).}

A Euclidean vector bundle has a Bott operator in every fibre. Orthogonal
changes of frame preserve its Clifford formula, so these operators form
a global cycle. The fibrewise Dirac operators form its inverse. A spinor
bundle then removes the Clifford coefficient algebra and gives the Thom
isomorphism. We prove the two inverse identities before making that
Morita reduction; this separates the analytic theorem from the choice of
K-orientation.

We use the preceding bounded connection, technical-partition and
associativity proofs, and the Clifford and Morita conventions in
\emph{Graded C*-algebras, Clifford algebras and graded Hilbert modules}.
Section 0 proves the bounded-transform and unbounded-product arguments
needed here. Sections 2--4 then prove the oscillator and
simultaneous-rotation inverse identities for the actual bundle cycles.
The relative-bundle description used below is proved in the earlier
\href{supporting/relative-k-foundations/relative-k-foundations.html\#rk-difference}{Relative
difference bundles, radial pairs and the planar normalization, Theorems
RK.3--RK.4}.
\href{supporting/relative-k-foundations/relative-k-foundations.html\#rk-planar-bott}{RK.6--RK.6a}
prove the actual planar projection and its complex-orientation sign. The
exact even relative character, splitting injection and ordinary
cohomological Thom inputs used in Section 8 are proved in
\href{supporting/cohomological-foundations/cohomological-foundations.html\#cf-thom}{Cohomological
foundations for the complex Thom comparison, CF.1--CF.6}.

For compact groups, the Haar probability measure used below is proved in
\href{supporting/representations-of-compact-groups/compact-foundations.html\#cpt-f-005}{Foundations
for compact-group averaging and coefficient approximation, CPT-F-005}.
Its
\href{supporting/representations-of-compact-groups/compact-foundations.html\#cpt-f-010}{CPT-F-010}
constructs continuous vector integrals, strict operator averages and
source-localized compact integrals on arbitrary Hilbert modules.

\subsection{0. The analytic product criterion used
below}\label{the-analytic-product-criterion-used-below}

We prove the analytic recognition argument here. Lemma 0.0 proves the
required resolvent, continuous functional-calculus and tensor-extension
foundations locally. The bounded product, technical partition and
essential-replacement proofs are the preceding programme lessons
\emph{Kasparov's technical theorem} and \emph{Connections and the
existence of the Kasparov product}. The following argument also treats a
nonessential first representation.

\textbf{Lemma 0.0 (resolvents, continuous calculus and tensor
extension).} Let \(E\) be a Hilbert \(B\)-module, and let \(D=D^*\) be a
closed densely defined operator for which \(1+D^2\) has dense range.
Then \(D\pm i\) are bijective from \(\operatorname{Dom}D\) to \(E\),
with mutually adjoint inverse resolvents of norm at most one. There is a
nondegenerate continuous functional calculus
\(C_0(\mathbb R)\to\mathcal L(E)\); it extends to bounded continuous
functions, and \[
\|(D-z)^{-1}\|\leq |\operatorname{Im}z|^{-1}
\quad(\operatorname{Im}z\ne0).
\] If \(Y\) is any \(B\)-\(C\) correspondence, the closure of
\(D\otimes1\) on \(\operatorname{Dom}D\odot_B Y\) is self-adjoint
regular on \(E\otimes_B Y\), and its bounded continuous functional
calculus is \(f(D)\otimes1\). No separability or essentiality of the
left action on \(Y\) is assumed.

\textbf{Proof: resolvents.} Symmetry cancels the mixed terms in \[
\langle(D\pm i)x,(D\pm i)x\rangle
=\langle Dx,Dx\rangle+\langle x,x\rangle.
\] Thus these operators are bounded below by one. Their ranges are
closed: a convergent sequence of images makes the preimages Cauchy, and
closedness of \(D\) then gives a preimage of the limit. The range of
each contains the dense range of \(1+D^2\), since on
\(\operatorname{Dom}D^2\) this is the product \((D+i)(D-i)\) in either
order. Their ranges are therefore all of \(E\). Write
\(R_\pm=(D\pm i)^{-1}\). For \(a,b\in E\), symmetry gives \[
\langle R_+a,b\rangle
=\langle R_+a,(D-i)R_-b\rangle
=\langle a,R_-b\rangle.
\] Hence \(R_+^*=R_-\). The identities \(DR_\pm=1\mp iR_\pm\), first on
their ranges, give \[
R_+-R_-=-2iR_+R_-=-2iR_-R_+.
\] In particular they commute. Put \(U=1-2iR_+\). The preceding
identities give \(UU^*=U^*U=1\), and \(1-U=2iR_+\) has dense range
\(\operatorname{Dom}D\).

\textbf{Proof: functional calculus.} We first prove the continuous
calculus of a unitary locally. For a Laurent polynomial
\(p(z)=\sum_{j=-m}^m a_jz^j\), put \(M=\sup_{|z|=1}|p(z)|\). Given
\(\xi\in E\) and \(N>2m\), the continuous \(E\)-valued polynomial \[
w_N(z)=N^{-1/2}\sum_{k=0}^{N-1}z^kU^{-k}\xi
\] has
\(\int_{\mathbb T}\langle w_N,w_N\rangle\,d\theta/(2\pi)=\langle\xi,\xi\rangle\).
For \(m\leq l\leq N-1-m\), the coefficient of \(z^l\) in \(p(z)w_N(z)\)
is \(N^{-1/2}U^{-l}p(U)\xi\). Integrating the inner product cancels
distinct monomials, and every omitted coefficient contributes a positive
element. Since pointwise \(|p(z)|^2\leq M^2\), this gives the
positive-order inequality \[
\frac{N-2m}{N}\langle p(U)\xi,p(U)\xi\rangle
\leq \int_{\mathbb T}\langle p w_N,p w_N\rangle\,\frac{d\theta}{2\pi}
\leq M^2\langle\xi,\xi\rangle.
\] Take norms and then \(N\to\infty\). Thus
\(\|p(U)\|\leq\|p\|_\infty\), for every Laurent polynomial, on a Hilbert
module as well as a Hilbert space.

Laurent polynomials are uniformly dense in \(C(\mathbb T)\). Indeed the
kernels \(K_N(\theta)=N^{-1}|\sum_{k=0}^{N-1}e^{ik\theta}|^2\) are
positive with normalized integral one. The geometric-series formula
bounds their integral outside any fixed neighborhood of zero by a
constant times \(N^{-1}\). Uniform continuity then makes convolution
with \(K_N\) converge uniformly to a continuous circle function; each
convolution is a Laurent polynomial. The polynomial norm bound therefore
extends evaluation at \(U\) to a unital contractive *-homomorphism
\(C(\mathbb T)\to\mathcal L(E)\). Multiplication and adjoints pass
through uniform polynomial approximation. No earlier normal-calculus or
spectral-representation assertion is needed for this construction. The
map \[
u(t)=(t-i)/(t+i)
\] identifies \(\mathbb R\) with the circle minus \(1\). For
\(f\in C_0(\mathbb R)\), its transported function extends by zero at
\(1\); evaluate it at \(U\) to define \(f(D)\). This is a contractive
*-homomorphism. It is nondegenerate because the function
\(1-u(t)=2i/(t+i)\) has operator value \(1-U\), whose range is dense.

For completeness, its extension to \(C_b(\mathbb R)\) is also
constructive. Let \(e_n\in C_c(\mathbb R)\) be positive cutoffs tending
uniformly to one on compact sets, with \(0\leq e_n\leq1\). For bounded
continuous \(h\), the uniformly bounded operators \((he_n)(D)\) converge
on every vector \(f(D)\xi\), because \(he_nf\to hf\) uniformly. Such
vectors have dense span, so the operators converge on every vector to a
bounded operator \(h(D)\). The same argument for \(\overline h\) gives
its adjoint. Products and adjoints pass to these uniformly bounded
vector limits, proving the extended *-homomorphism. Its value does not
depend on the chosen cutoffs, by the same dense-span argument.

The identity \(D=i(1+U)(1-U)^{-1}\) holds on the dense domain
\(\operatorname{ran}(1-U)\). For nonreal \(z\), the bounded scalar
function \((t+i)/(t-z)\) therefore shows that the value of
\((t-z)^{-1}\) has range in \(\operatorname{Dom}D\) and is a two-sided
inverse of \(D-z\). Its norm is bounded by the scalar supremum
\(|\operatorname{Im}z|^{-1}\). Polynomial identities on \(U\) and their
uniform limits show that this calculus agrees with the two original
resolvents and with the usual products of their functions. In particular
the value of \((1+t^2)^{-1}\) is \((1+D^2)^{-1}=R_+R_-\).

We also need the converse Cayley construction. If \(V\) is unitary and
\(1-V\) has dense range, then it is injective: its kernel equals the
kernel of \(1-V^*\), which is orthogonal to that dense range. On
\(\operatorname{ran}(1-V)\) define \[
S((1-V)x)=i(1+V)x.
\] Expanding both inner products and using \(V^*V=1\) proves symmetry.
The operators \(S+i\) and \(S-i\) have bounded, everywhere-defined,
mutually adjoint inverses \[
Q_+=(1-V)/(2i),\qquad Q_-=(V^*-1)/(2i).
\] These formulas prove that \(S\) is closed, since the inverse of an
injective bounded operator is closed on its range. If
\(y\in\operatorname{Dom}S^*\), choose \(x\in\operatorname{Dom}S\) with
\((S-i)x=(S^*-i)y\). Then \(y-x\in\ker(S^*-i)\), which is orthogonal to
the surjective range of \(S+i\). Hence \(y=x\), proving \(S=S^*\). The
commuting product \(Q_+Q_-\) has range in \(\operatorname{Dom}S^2\) and
is the inverse of \(1+S^2\), by the two resolvent identities. Thus \(S\)
is regular. This proves the converse without a graph-complement theorem.

\textbf{Proof: tensor extension.} Lemma 6.0a of Lesson 05 supplies the
unital *-homomorphism \(T\mapsto T\otimes1\) on adjointable operators.
Thus \(V=U\otimes1\) is unitary. The range of \(1-V\) is dense: the
algebraic tensors with first vector in \(\operatorname{Dom}D\) are
dense, and every such vector is in that range. Apply the converse Cayley
construction to \(V\), obtaining a self-adjoint regular \(S\). The
formulas for \(Q_\pm\) show that \(S(x\otimes y)=Dx\otimes y\) on the
algebraic domain. This domain is a graph core: its image under \(S+i\)
contains all elementary tensors and is dense, while \(S+i\), with its
bounded inverse and \(S Q_+=1-iQ_+\), identifies its graph norm with an
equivalent complete norm on \(E\otimes_B Y\). Consequently \(S\) is
precisely the closure of the algebraic tensor operator. Continuous
circle calculus commutes with \(T\mapsto T\otimes1\). The
bounded-function extension does too, since uniformly bounded vector
convergence on \(E\) gives vector convergence on elementary tensors and
hence on their dense span. This proves the claimed calculus identity. If
\(D\) is odd, its grading sends \(U\) to \(U^*\), and the same formulas
give \(f(D)\mapsto f(-D)\) and oddness of the tensor extension.
\(\square\)

The free
\href{https://www.math.uni-bonn.de/people/lesch/dl/pap/2011-KaaLes-ArXiv-v2.pdf}{Kaad--Lesch
author preprint, Section 2.3 and Proposition 4.1} supplies the
regular-operator formulation used here. Its references are not imported
as proof premises: the resolvent, Cayley, functional-calculus and tensor
assertions needed below have just been proved.

An unbounded module \((E,\phi,D)\) means a countably generated graded
Hilbert module, a graded representation, and an odd self-adjoint regular
operator. We require \((1+D^2)^{-1}\phi(a)\) compact for every source
element. A dense graded \(*\)-subalgebra preserves
\(\operatorname{Dom}D\) and has bounded adjointable graded commutators
with \(D\). The representation need not be essential.

\textbf{Lemma 0.1 (bounded transformation with local compactness).} For
an unbounded module, \(F_D=D(1+D^2)^{-1/2}\) is a Kasparov-cycle
operator.

\textbf{Proof.} Put \(q(t)=(1+t^2)^{-1}\). Cutoffs \(e_n\) on \((0,1]\),
zero near zero and eventually one on every compact subinterval, make
\(e_n(q)/q\) bounded. Hence \(e_n(q(D))\phi(a)\) is compact. For
\(f\in C_0(\mathbb R)\), the functions \(fe_n(q)\) converge uniformly to
\(f\); therefore \(f(D)\phi(a)\) is compact. Adjointing gives
compactness on the other side.

For a homogeneous smooth source element write \(\epsilon=(-1)^{|a|}\),
\(c=[D,\phi(a)]_{\rm gr}\), and \(R_z=(D-z)^{-1}\). Domain preservation,
followed by multiplication by inverse resolvents, gives \[
R_z\phi(a)-\epsilon\phi(a)R_{\epsilon z}
=-\epsilon R_zcR_{\epsilon z},\qquad z=\pm i\rho.
\tag{0.1}
\] For \(Q_\rho=D(D^2+\rho^2)^{-1}=(R_{i\rho}+R_{-i\rho})/2\), it
follows that \(\|[Q_\rho,\phi(a)]_{\rm gr}\|\leq\|c\|\rho^{-2}\). The
commutator becomes compact after source multiplication on either side,
since the corresponding end resolvent is locally compact.

Functional calculus gives the vector integral \[
F_D\xi=\frac2\pi\int_0^\infty
D(1+D^2+t^2)^{-1}\xi\,dt.
\tag{0.2}
\] Indeed its partial-integral scalar functions are bounded by one,
converge uniformly on compact subsets of the real line, and converge
pointwise to \(x/\sqrt{1+x^2}\). Apply them first to the dense set of
vectors \(f(D)\eta\), \(f\in C_0(\mathbb R)\), and then use the uniform
bound. The uncommuted integral need not converge in operator norm. Its
commutator does: (0.1) bounds its integrand by \(\|c\|/(1+t^2)\). Thus
both source-localized commutators of \(F_D\) are compact. For
homogeneous smooth \(a,b\), the identity \[
[F_D,\phi(ab)]_{\rm gr}
=[F_D,\phi(a)]_{\rm gr}\phi(b)
 +(-1)^{|a|}\phi(a)[F_D,\phi(b)]_{\rm gr}
\] makes the unlocalized commutator for \(ab\) compact. Products from
the dense smooth algebra linearly span a dense subset of the source:
approximate an algebra approximate identity and the given element by
smooth elements and multiply. Since the commutator is bounded by
\(2\|a\|\) on each parity subspace, density proves compactness for every
source element. Finally \(F_D\) is odd and self-adjoint, and
\(F_D^2-1=-(1+D^2)^{-1}\). These are all cycle conditions. \(\square\)

Write \(E=E_1\widehat\otimes_B E_2\), \(\phi=\phi_1\widehat\otimes1\),
and \(E_1^0=\overline{\phi_1(A)E_1}\). For a homogeneous vector \(x\),
\(T_x\eta=x\widehat\otimes\eta\) is the creation operator.

\textbf{Lemma 0.2 (the connection estimate).} Suppose a homogeneous
linear subspace \(\mathcal X\subset\operatorname{Dom}D_1\), dense in
\(E_1^0\), satisfies the creation and adjoint domain inclusions and has
bounded adjointable errors \[
DT_x-(-1)^{|x|}T_xD_2,\qquad
T_x^*D-(-1)^{|x|}D_2T_x^*.
\tag{0.3}
\] Then \(F_D\) satisfies both \(F_{D_2}\)-connection conditions for
every \(x\in E_1^0\).

\textbf{Proof.} On \(E\oplus E_2\), the two creations are the
off-diagonal blocks of a homogeneous adjointable operator. Their domain
inclusions and (0.3) give a bounded graded commutator with
\(D\oplus D_2\). The inverse-resolvent identity (0.1) applies to this
operator too. Multiply its bounded-transform commutator on the right by
\(\operatorname{diag}(0,\phi_2(b))\), for a homogeneous smooth
second-source element \(b\). The localized end resolvent is compact, and
the bound \(C/(1+t^2)\) in (0.2) is integrable. The top-right block
proves \[
\bigl(F_DT_x-(-1)^{|x|}T_xF_{D_2}\bigr)\phi_2(b)
\] compact. Since \(T_{xb}=T_x\phi_2(b)\), moving \(F_{D_2}\) past
\(\phi_2(b)\) adds only the compact commutator proved in Lemma 0.1. Thus
the creation error for \(xb\) is compact; adjointing gives its second
condition. The linear span of these vectors is dense in \(E_1^0\), by
nondegeneracy of its right coefficient action, density of
\(\mathcal X\), and density of the smooth coefficient algebra. The
inequality \(\|T_x\|\leq\|x\|\) extends both assertions to every vector
of \(E_1^0\). \(\square\)

\textbf{Lemma 0.3 (bounded alignment, including essential connections).}
Let \(A\) be separable and the intermediate algebra be
\(\sigma\)-unital. Suppose \(F_1\) is a normalized first-cycle operator
and \(G\) is a self-adjoint contraction cycle on \(E\), with the
\(F_2\)-connection conditions on \(E_1^0\). If, for a fixed
\(0\leq\kappa<2\), \[
\phi(a)[F_1\widehat\otimes1,G]_{\rm gr}\phi(a)^*
\geq-\kappa\phi(aa^*)\pmod{\mathcal K(E)}
\qquad(a\in A),
\tag{0.4}
\] then \(G\) represents the Kasparov product.

\textbf{Proof.} First assume a full connection. The product-existence
proof, using this connection and its technical partition \(M+N=1\),
gives a product operator \(H=M^{1/4}SM^{1/4}+N^{1/4}GN^{1/4}\),
\(S=F_1\widehat\otimes1\). The same root commutation and defect
annihilation used there give, after a source sandwich, \[
[G,H]_{\rm gr}
\equiv M^{1/2}[G,S]_{\rm gr}+2N^{1/2}
\pmod{\text{source-localized compacts}}.
\] Consequently (0.4), \(0\leq M^{1/2}\leq1\), and positivity of the
second term bound this sandwich below by \(-\kappa\phi(aa^*)\).

Here are the quotient and homotopy details of that conclusion. Let
\(\mathcal D_\phi\) be the graded algebra of bounded operators whose
graded source commutators are compact, and let \(\mathcal C_\phi\) be
its ideal of two-sided source-locally compact operators. Closure under
products and adjoints follows from the graded Leibniz identity. For an
even self-adjoint \(Z\in\mathcal D_\phi\), positivity of every source
sandwich modulo compacts implies \(Z_-\in\mathcal C_\phi\). To see this,
work modulo \(\mathcal K(E)\), where \(Z\) commutes with the represented
source. The negative part of \(\phi(aa^*)Z\) is \(\phi(aa^*)Z_-\); it is
zero. Hence \(Z_-\phi(a)\) has zero product with its adjoint in this
quotient and is compact. Adjointing gives the other side. Continuous
functional calculus is legitimate because it follows by polynomial
approximation in the unital algebra generated by \(Z\).

Apply this to \(Z=[G,H]_{\rm gr}+\kappa\). In
\(\mathcal D_\phi/\mathcal C_\phi\), the cycle operators \(G,H\) are
self-adjoint unitaries and their anticommutator is at least \(-\kappa\).
For \(L_t=(1-t)G+tH\) this gives \[
L_t^2\geq (1-t)^2+t^2-\kappa t(1-t)
\geq(2-\kappa)/4.
\] Choose \(0<\delta<(2-\kappa)/4\) and normalize by
\(L_t\bigl(\max(L_t^2,\delta)\bigr)^{-1/2}\). This is a norm-continuous
path of cycles: it is an odd self-adjoint unitary in the defect
quotient, and its source commutators are compact by functional calculus.
At its endpoints normalization changes the original cycle only by a
two-sided locally compact operator. The straight interpolation for such
a change is a cycle, as follows by expanding its square and source
commutator. Thus \(G\) and \(H\) have the same class.

We now remove the full-connection assumption. Put
\(E^0=\overline{\phi(A)E}\); the balancing rule identifies it with
\(E_1^0\widehat\otimes_B E_2\). Proposition 5.1 of the product lesson
supplies essential replacement cycles \(F_1^0,G^0\) on these modules in
the original classes. Its connection relations for multiplication maps
\(A_a^1:E_1\to E_1^0\), \(A_a:E\to E^0\) are \[
F_1^0A_a^1-(-1)^{|a|}A_a^1F_1\in\mathcal K(E_1,E_1^0),\qquad
G^0A_a-(-1)^{|a|}A_aG\in\mathcal K(E,E^0).
\tag{0.5}
\] These maps are adjointable: their adjoints are multiplication by
\(a^*\) followed by inclusion. Normalizing the replacement cycles
preserves (0.5).

The operator \(G^0\) is a full connection on \(E_1^0\). Indeed
\(T_{ax}=A_aT_x\). Expanding its error with (0.5) leaves a compact term
and the source-localized creation error of \(G\). That latter term is
compact after inclusion in \(E\), because the connection error for
\(ax\in E_1^0\) is compact and the source commutator of \(G\) is
compact. It is compact with target \(E^0\) as well: for an adjointable
map \(R\), compactness of \(R^*R\) implies compactness of \(R\). In
fact, with \(q_\varepsilon(t)=t/(t+\varepsilon)\), the maps
\(Rq_\varepsilon(R^*R)\) are compact and converge in norm to \(R\),
since \(\sup_{t\geq0}\sqrt t\,\varepsilon/(t+\varepsilon)\to0\). Taking
adjoints and using density of the vectors \(ax\) proves both full
connection conditions.

For \(k\in\mathcal K(E_1)\), the weak connection also gives \[
\phi(a)[G,k\widehat\otimes1]_{\rm gr}\phi(b)\in\mathcal K(E).
\tag{0.6}
\] Indeed \(\phi_1(a)k\phi_1(b)\) is a norm limit of rank-one operators
with both vectors in \(E_1^0\). Their tensor images are \(T_xT_y^*\);
the connection identities cancel the two \(F_2\) terms in their
commutator. Expanding the commutator with the two outer source factors
adds only compact source commutators of \(G\). This proves (0.6).

Take even \(c\in A\). By (0.5),
\(R_c=(A_c^1)^*F_1^0A_c^1-\phi_1(c)^*F_1\phi_1(c)\) is compact. Put
\(S^0=F_1^0\widehat\otimes1\). Moving \(G^0\) past \(A_c\), then
sandwiching by \(\phi(d)\), gives \[
\phi(d)^*A_c^*[S^0,G^0]_{\rm gr}A_c\phi(d)
\equiv\phi(cd)^*[S,G]_{\rm gr}\phi(cd)
\pmod{\mathcal K(E)}.
\tag{0.7}
\] The extra bracket of \(R_c\widehat\otimes1\) is compact by (0.6); the
remaining errors are (0.5) and compact source commutators. Equation
(0.4), applied with \(a=(cd)^*\), transfers its lower bound through
(0.7).

This transfer holds on \(E^0\), not only after inclusion into \(E\). A
compact operator on \(E\) whose range and adjoint range lie in \(E^0\)
belongs to \(\mathcal K(E^0)\). For a proof, first regard it as an
adjointable map \(E\to E^0\); the compactness test just proved makes
that map compact. An approximate identity \(u_\lambda\) of
\(\mathcal K(E^0)\) then has \(u_\lambda T\to T\) in norm, by finite
rank approximation with target vectors in \(E^0\). Apply the same
argument to \(T^*\) to obtain \(Tu_\lambda\to T\). The sandwiches
\(u_\lambda T u_\lambda\) lie in \(\mathcal K(E^0)\) and converge to
\(T\). For the source sandwich in (0.7), its negative part has the same
range property, since its continuous defining function vanishes at zero.
Thus its compact negative part restricts compactly to \(E^0\). Positive
even approximate identities \(c\) give \(cd\to d\), proving (0.4) for
all sources on the essential module. The full-connection case applies
there. Essential replacement and homotopy invariance of products return
the assertion to the original modules. \(\square\)

\textbf{Lemma 0.4 (integrated positivity).} Let \(D,S\) be odd
self-adjoint regular operators with
\(\operatorname{Dom}D\subset\operatorname{Dom}S\). If \[
Q(\xi):=2\operatorname{Re}\langle D\xi,S\xi\rangle
\geq-c\langle\xi,\xi\rangle\quad(\xi\in\operatorname{Dom}D),
\tag{0.8}
\] for \(c\geq0\), then \([F_D,F_S]_{\rm gr}\geq-c\). Replacing \(S\) by
\(\alpha S\), \(\alpha>0\), replaces this bound by \(-\alpha c\).

\textbf{Proof.} For \(\lambda=1+u^2\), \(\mu=1+v^2\), set \[
\begin{gathered}
r_D=(D^2+\lambda)^{-1},\quad h_D=Dr_D,\quad k_D=\sqrt\lambda r_D,\\
r_S=(S^2+\mu)^{-1},\quad h_S=Sr_S,\quad k_S=\sqrt\mu r_S.
\end{gathered}
\] For \(\eta\in\operatorname{Dom}S\), expanding \(Dh_D=1-\lambda r_D\),
\(Dk_D=\sqrt\lambda h_D\) gives \[
Q(h_D\eta)+Q(k_D\eta)=2\operatorname{Re}\langle S\eta,h_D\eta\rangle.
\] The cancelling difference is
\(2\lambda\operatorname{Re}(\langle h_D\eta,Sr_D\eta\rangle-\langle r_D\eta,Sh_D\eta\rangle)\),
zero by symmetry of \(S\). All these vectors are in the indicated
domains by the domain inclusion. Apply the identity to
\(h_S\psi,k_S\psi\). Using \(Sh_S=1-\mu r_S\), \(Sk_S=\sqrt\mu h_S\),
the remaining two terms cancel by self-adjointness of \(h_D\). Hence \[
\sum_{w\in\{h_Dh_S,k_Dh_S,h_Dk_S,k_Dk_S\}}Q(w\psi)
=2\operatorname{Re}\langle\psi,h_Dh_S\psi\rangle.
\tag{0.9}
\] Integrate first over finite rectangles \(0\leq u,v\leq R\), with
factor \(4/\pi^2\). Formula (0.2) makes the right side converge to the
quadratic form of the bounded anticommutator. The lower bound (0.8)
bounds the left side below by minus \(c\) times the sum of the four
positive forms \(\langle w\psi,w\psi\rangle\). Their total integral is
at most \(\langle\psi,\psi\rangle\): first use \[
h_D^2+k_D^2=r_D,\qquad
\frac2\pi\int_0^\infty r_D\,du=(1+D^2)^{-1/2}\leq1,
\] then bound the two remaining sandwiches by \(h_S^2+k_S^2=r_S\) and
integrate once more. These increasing positive integrals also justify
passage to the rectangle limit without asserting absolute integrability
of the unsplit anticommutator integrand. The inequality holds for every
\(\psi\), hence as an operator inequality. Scaling \(S\) scales (0.8) by
\(\alpha\), proving the last assertion. \(\square\)

\textbf{Theorem 0.5 (the unbounded product criterion).} Let \(A\) be
separable and \(B\) be \(\sigma\)-unital. Suppose \((E_1,\phi_1,D_1)\),
\((E_2,\phi_2,D_2)\), and
\((E_1\widehat\otimes_B E_2,\phi_1\widehat\otimes1,D)\) are unbounded
modules. Suppose (0.3) holds on a homogeneous subspace of
\(\operatorname{Dom}D_1\) dense in \(\overline{\phi_1(A)E_1}\), with
both creation domain inclusions. Suppose
\(\operatorname{Dom}D\subset\operatorname{Dom}(D_1\widehat\otimes1)\),
and (0.8) holds with \(S=D_1\widehat\otimes1\). Then the
bounded-transform class of \(D\) is the Kasparov product of the classes
of \(D_1,D_2\). The first representation may be nonessential.

\textbf{Proof.} Lemma 0.1 supplies the cycles and Lemma 0.2 their
essential-source connection. Choose \(\alpha>0\) with \(\alpha c<2\).
Lemma 0.4 gives (0.4) for \(F_D\) and
\(F_{\alpha D_1}\widehat\otimes1\); functional calculus commutes with
the tensor representation. Lemma 0.3 identifies their product. Finally
positive rescaling preserves the class of the first operator: on a
compact positive interval of parameters, \(tx/\sqrt{1+t^2x^2}\) varies
uniformly in \(x\), continuously in \(t\). Its transforms form a
norm-continuous path of cycles by Lemma 0.1. Domains and bounded
commutators are unchanged by this positive scaling, and local
compactness follows from the vanishing-functions argument in its proof.
Thus \([\alpha D_1]=[D_1]\), proving the assertion. \(\square\)

The freely readable \href{https://arxiv.org/pdf/2006.10616v2}{van den
Dungen preprint, Sections 2.4 and 3.1}, printed pp.~11--16, provides the
comparison account for connection estimates and integrated positivity.
The proof used here is the complete local argument above. In particular
its nonessential representation step does not rely on the preprint's
reference to an external original proof.

\subsection{1. Central coefficients and compact
bases}\label{central-coefficients-and-compact-bases}

Let \(X\) be compact Hausdorff and put \(R=C(X)\), trivially graded. An
\(R\)-algebra is a graded C*-algebra with a specified even central
action of \(R\). An \(R\)-linear cycle has the property that the left
and right actions of \(R\) on its module agree. Write \(KK_R(A,B)\) for
its homotopy classes. Forgetting this agreement gives a usual KK class.

The tensor products of modules below are balanced over their displayed
coefficient algebra. External products over \(R\) use the balanced
algebra tensor product over \(R\). Thus both factors refer to the same
point of \(X\).

\textbf{Lemma 1.1 (the central coefficient argument).} Suppose the
closed \(R\)-linear *-algebra generated by countably many elements is
\(A\), and all coefficient and intermediate algebras under consideration
are \(\sigma\)-unital. For \(R\)-linear countably generated cycles, the
product-existence, uniqueness, associativity and unbounded
product-recognition arguments used earlier apply with this hypothesis in
place of separability of \(A\). These operations preserve
\(R\)-linearity.

\textbf{Proof.} We specify the modification, since \(C(X)\) need not be
separable. Every adjointable right-module operator commutes with the
right action of the central algebra \(R\), hence with its left action in
an \(R\)-linear cycle. This includes all compact operators, connection
operators, technical partitions and their functional calculi.

In the technical theorem, take as the derivating space the separable
space generated by the countable noncentral source generators, the
required cycle and connection operators, and their adjoints. Include
countably many words in the source generators. All the ideals used in
the product proof remain \(\sigma\)-unital: the compact ideal of a
countably generated module is \(\sigma\)-unital, its represented images
are \(\sigma\)-unital, and the proof uses only sums of these images and
finitely or countably many added operators. The technical theorem
supplies the same partition. Its commutators with the omitted elements
of \(R\) are exactly zero.

The source conditions extend from the selected words to products with
elements of \(R\), since the latter commute with every operator in the
argument. Finite sums of these products are dense in \(A\). Norm
continuity then gives the cycle, connection and localized positivity
conditions for every source element. This proves existence and the
uniqueness comparison by the same convex-path and normalization argument
as before.

For associativity, use precisely the two compact images and the
localized negative-part ideal in Theorem 4.1 of \emph{Homotopy,
associativity, the index pairing and KK-equivalence}. The extra
derivating operators there are again countable; central coefficients
commute with them and with the partition. Both positivity terms in that
proof are retained. This gives associativity whenever the displayed
binary products exist. The argument for balanced external products is
the identical argument on the canonically associated module tensor
products.

Finally the creation-operator resolvent identities and the integrated
positivity estimate in the unbounded product criterion require no
separable source. Separability enters its product comparison only
through the technical partition just described. Thus its domain,
connection and lower-bound hypotheses give product recognition here as
well. All operators constructed are right-module operators, so remain
\(R\)-linear. The same construction on interval modules proves
compatibility with homotopies. \(\square\)

This lemma concerns central module actions. It makes no assertion about
an arbitrary representation of a nonseparable source algebra.

Let \(V\to X\) be a rank-\(n\) Euclidean real vector bundle. Write \[
\begin{aligned}
C&=\Gamma(\operatorname{Cl}(V)),\\
B&=C_0(V),\\
A&=B\widehat\otimes_R C.
\end{aligned}
\tag{1.1}
\] The complex Clifford convention is \(vw+wv=2(v,w)\). The Clifford
algebra is graded; \(B\) is not.

The algebra \(C\) is unital. The radial functions of \(\|v\|\), equal to
one on expanding balls, give a countable approximate identity of \(B\)
and \(A\). All three are generated over \(R\) by countably many
elements. For \(C\), take finitely many bundle sections spanning each
fibre, obtained from a finite trivializing cover and a partition of
unity. For \(B\), embed \(V\) isometrically in a finite trivial bundle
and use its coordinate functions multiplied by countably many radial
cutoffs. Polynomials in these coordinates approximate continuous
functions on bounded fibre balls; a finite bundle cover and a partition
of unity make the approximation uniform over \(X\). Cut off at larger
radii to approximate \(C_0(V)\). The same argument with the finite
Clifford sections proves the assertion for \(A\). Thus Lemma 1.1 applies
without a metrizability hypothesis on \(X\).

Rank may be locally constant rather than constant. A finite-rank bundle
on a compact base has only finitely many ranks; their inverse images are
clopen. Apply the constructions separately and take the finite direct
sum. We write constant rank to simplify the formulas.

\subsection{2. The fibrewise position and Dirac
cycles}\label{the-fibrewise-position-and-dirac-cycles}

On the standard right \(A\)-module let \[
\begin{gathered}
X_V(v)=v\in\operatorname{Cl}(V_{\pi(v)}),\\
F_V(v)=\frac{v}{\sqrt{1+\|v\|^2}}.
\end{gathered}
\tag{2.1}
\] The source \(R\) acts by pullback along \(\pi:V\to X\).

\textbf{Proposition 2.1 (the position class).} Formula (2.1) defines \[
b_V\in KK_R(R,A).
\tag{2.2}
\]

\textbf{Proof.} It is odd and self-adjoint, its source commutators
vanish, and \[
F_V^2-1=-\frac1{1+\|v\|^2}\in A.
\tag{2.3}
\] This is a function vanishing at infinity because the base is compact
and the norm-bounded ball bundles are compact. Multiplication by it is
compact on the standard module. The module is countably generated by the
radial approximate identity. These verify every cycle condition. The
unbounded multiplier also has the explicit inverse resolvents \[
(X_V\pm i)^{-1}=\frac{X_V\mp i}{1+\|v\|^2};
\] their dense ranges contain the compactly supported continuous
sections. Thus it is self-adjoint and regular. \(\square\)

Over \(x\in X\) put \[
\mathcal H_x=L^2(V_x)\otimes\Lambda^*(V_x\otimes_{\mathbb R}\mathbb C).
\tag{2.4}
\] Orthogonal frame changes act by their unitary change of variable on
\(L^2\) and their exterior action on forms. These give a continuous
Hilbert bundle: its transition operators are strongly continuous, which
suffices for continuity of its sections. Let \(H_V\) be the Hilbert
\(R\)-module of continuous sections. Its inner product is the fibre
integral.

In an orthonormal frame set \[
\begin{gathered}
c_j=\varepsilon_j+\iota_j,\\
\widehat c_j=\varepsilon_j-\iota_j,\\
D_V=\sum_j\widehat c_j\partial_j.
\end{gathered}
\tag{2.5}
\] Both Clifford families transform with the frame, so the operator and
the representation of \(A\) by multiplication and \(c_j\) are globally
defined. Grade \(H_V\) by form parity.

\textbf{Proposition 2.2 (the vertical Dirac class).} The fibre operator
(2.5) defines an odd self-adjoint regular operator on \(H_V\). Its
domain is the continuous field of first Sobolev domains with the graph
norm. Its bounded transform gives \[
a_V\in KK_R(A,R).
\tag{2.6}
\]

\textbf{Proof.} In a bundle chart this is the fixed Euclidean operator
from Lemma 3.1 of the Bott lesson. Its inverse resolvents are uniformly
bounded, and are carried into one another by the frame-transition
unitaries. They therefore give adjointable global operators, with
adjoints the opposite resolvents. Their ranges contain sections locally
given by compactly supported smooth fibre vectors, times continuous base
functions; finite partitions of unity show that such sections are dense.
The resolvent identities, symmetry and dense range prove
self-adjointness and regularity on the stated graph domain.

Countable generation also follows from a finite cover. Choose partitions
with supports in bundle charts and a countable dense family of smooth
fibre vectors in each chart. Multiplying these by the partition
functions gives countably many global sections; their \(R\)-linear span
is dense. Strong continuity of the transition unitaries is enough for
this assertion.

For local compactness take a section of \(A\) supported in one chart and
compactly supported in its fibre variable. Approximate it uniformly by
finite sums of a continuous base coefficient times a compactly supported
fibre function and a Clifford matrix. The localized Euclidean inverse
resolvent is compact by the frequency-cutoff and
square-integrable-kernel proof in the Bott lesson. Multiplying its
finite rank approximations on each side by chart partition functions
gives finite rank operators on \(H_V\). Summing finitely many chart
pieces and taking norm limits proves local compactness for every element
of \(A\).

There is a dense graded *-algebra of sections continuous in the base and
smooth with compact support in the fibre in each chart. On it the graded
commutator is Clifford multiplication by the fibre derivative, a bounded
multiplier. Domain preservation is the fibre Sobolev product rule. The
bounded-transform theorem now applies. All operators commute with \(R\),
so the cycle is \(R\)-linear. \(\square\)

\subsection{3. The two inverse
identities}\label{the-two-inverse-identities}

\textbf{Theorem 3.1 (the parametrized Bott theorem).} The classes above
satisfy \[
\begin{aligned}
b_V\widehat\otimes_A a_V&=1_R,\\
a_V\widehat\otimes_R b_V&=1_A.
\end{aligned}
\tag{3.1}
\]

\textbf{Proof of the first identity.} Identify the tensor module
\(A\widehat\otimes_A H_V\) with \(H_V\). The candidate product is the
bounded transform of \[
Q_V=D_V+\sum_jv_jc_j.
\tag{3.2}
\] Its graph domain is the field of weighted Sobolev spaces \[
\{u:\partial_j u,\ v_j u\in L^2\text{ for every }j\}.
\tag{3.3}
\] The earlier local oscillator proof in \emph{Bott periodicity in KK:
the Bott and Dirac elements}, Lemma 4.0, proves the Euclidean
self-adjoint closure, compact graph inclusion, even Gaussian kernel and
zero odd kernel on this domain. Its Lemma 0.0a proves the Euclidean
Fourier unitarity used in Proposition 2.2. These are frame-invariant
statements. The chart resolvents glue as in Proposition 2.2.

Here the inverse resolvent is globally compact on the Hilbert
\(R\)-module. Indeed the Euclidean oscillator inverse resolvent is
compact, so its conjugation by the strongly continuous frame unitaries
is norm continuous. This follows first for rank-one operators, then for
compact operators by uniform finite rank approximation. In finitely many
charts approximate this field in norm by finite rank operators, multiply
by partition functions on both sides, and sum. The resulting operators
are finite sums of global rank-one operators.

The creation operator associated to a smooth compactly supported
Clifford section \(d\) is multiplication by \(d\). Its graded error
relative to \(D_V\) is \[
\begin{gathered}
Q_V M_d-(-1)^{|d|}M_dD_V\\
 =[D_V,M_d]_{\rm gr}+X_VM_d,
\end{gathered}
\tag{3.4}
\] which is bounded. Its adjoint satisfies the corresponding domain and
error conditions. The domain (3.3) is contained in that of \(X_V\), and
on its smooth core \[
\begin{gathered}
\{Q_V,X_V\}\\
=2\|v\|^2+2N-n\geq -n.
\end{gathered}
\tag{3.5}
\] Here \(N\) is form degree. The inequality extends to the graph
domain. The unbounded product criterion, with the central coefficient
argument of Lemma 1.1, identifies (3.2) as the product.

The normalized Gaussian \[
g_x(v)=\pi^{-n/4}e^{-\|v\|^2/2}
\tag{3.6}
\] is a global even unit section of the kernel, so that kernel is the
module \(R\). The complementary operator is invertible, with a uniform
spectral gap: every fibre is unitarily the same Euclidean oscillator
with its kernel removed. Deform its bounded transform to its sign,
leaving the kernel at zero. This is a norm-continuous
functional-calculus homotopy; its defects are compact and its source
commutators are zero. The complementary summand becomes degenerate. The
kernel summand is \((R,1,0)\), proving the first identity.

\textbf{Proof of the reverse identity.} Work on \(V\oplus V\), with
vectors \(v,w\) in the same fibre and Clifford families \(e_j,f_j\). For
\(0\leq\theta\leq\pi/2\) put \[
\begin{aligned}
u&=\cos\theta\,v+\sin\theta\,w,\\
t&=-\sin\theta\,v+\cos\theta\,w,\\
e_j^\theta&=\cos\theta\,e_j+\sin\theta\,f_j,\\
f_j^\theta&=-\sin\theta\,e_j+\cos\theta\,f_j .
\end{aligned}
\tag{3.7}
\] These formulas commute with every orthogonal change of frame, so
define global data. Use the interval standard module over
\(A\widehat\otimes_R A\), source \(A\) represented by its functions in
\(t\) and its Clifford generators \(f_j^\theta\), and operator \[
\frac{\sum_ju_je_j^\theta}{\sqrt{1+\|u\|^2}}.
\tag{3.8}
\] The graded source commutators vanish. Its localized square defect
vanishes uniformly at infinity: \(\|u\|^2+\|t\|^2=\|v\|^2+\|w\|^2\), so
either the inverse quadratic factor is small or the source function is
small. Joint continuity gives strict multiplier continuity and an actual
interval cycle. Operator-norm continuity of the multiplier is
unnecessary.

The endpoints prove \[
b_V\boxtimes_R1_A=[\rho]\boxtimes_R b_V,
\tag{3.9}
\] where \(\rho\) reflects both the fibre coordinate and the Clifford
vector. Compose with \(1_A\boxtimes_R a_V\). The left side is
\(a_Vb_V\); the right side is \([\rho]\), by the first identity. Thus
\(p=a_Vb_V\) is invertible, since \(\rho\) is an automorphism. It is
also idempotent: \[
p^2=a_V(b_Va_V)b_V=p.
\tag{3.10}
\] An invertible idempotent is the identity, proving the second equation
of (3.1). Every product and homotopy is \(R\)-linear; Lemma 1.1
justifies the displayed associativity even for nonmetrizable \(X\).
\(\square\)

\subsection{4. Removing the doubled Clifford
algebra}\label{removing-the-doubled-clifford-algebra}

There is a useful version with the Clifford algebra in the source rather
than the coefficient: \[
\begin{gathered}
x_V\in KK_R(C,B),\\
y_V\in KK_R(B,C).
\end{gathered}
\tag{4.1}
\] We describe its Morita correspondence, including its grading.

Let \(M=\Gamma(\Lambda^*V_{\mathbb C})\) as a finite Hilbert
\(R\)-module. Represent the first copy of \(C\) by \(c_j\), and the
second by \(i\widehat c_j\). These are odd self-adjoint anticommuting
Clifford families. Give \(M\) the grading \[
\Gamma_M=(-1)^{n(n-1)/2}(-1)^N.
\tag{4.2}
\] It is the chirality for the ordered doubled frame, first all \(e_j\),
then all \(f_j\). To check the sign, the chirality for the interleaved
frame is \(\prod_j(-i c_j\,i\widehat c_j)=(-1)^N\); sorting that frame
takes \(n(n-1)/2\) exchanges of odd generators.

\textbf{Lemma 4.1 (double Clifford Morita correspondence).} The
representation gives a graded isomorphism \[
C\widehat\otimes_R C\cong\operatorname{End}_R(M).
\tag{4.3}
\] Thus \(M\) is a graded Morita correspondence to \(R\).

\textbf{Proof.} In a fibre, the two Clifford families recover exterior
multiplication and contraction by \[
\varepsilon_j=(c_j-i f_j)/2,\qquad
\iota_j=(c_j+i f_j)/2.
\] The product \(\prod_j\iota_j\varepsilon_j\) is the vacuum projection.
Multiplication on the left and right by exterior creations and
contractions gives every matrix unit in the exterior basis. The
representation is onto. Both sides have dimension \(2^{2n}\), so it is
injective. The construction respects frame changes and (4.2), proving
the bundle and section assertion.

The usual left inner product is the rank-one operator
\(\theta_{\xi,\eta}\), identified through (4.3), and the right inner
product is the Hermitian fibre inner product. Their compatibility is the
rank-one formula. Finite partitions of unity and local frames make the
module finitely generated projective. Its inner products span both
coefficient algebras; the right span contains the constant one through
the vacuum section. Its dual is the inverse Morita correspondence by
evaluation of these inner products. All left actions are compact, so the
zero operators give the two Morita cycles and their explicit inverse
products. \(\square\)

Extend \(M\) to \(B\), denoting that correspondence by \(M_B\). Define
\[
\begin{aligned}
x_V&=(b_V\boxtimes_R1_C)
       \widehat\otimes_{B\widehat\otimes_R C\widehat\otimes_R C}M_B,\\
y_V&=M_B^*
       \widehat\otimes_{B\widehat\otimes_R C\widehat\otimes_R C}
          (a_V\boxtimes_R1_C).
\end{aligned}
\tag{4.4}
\] The first Clifford copy is the position copy in \(b_V\); the source
\(C\) is the second copy. This order matters.

\textbf{Theorem 4.2 (Clifford Thom equivalence).} For every compact
Hausdorff base, \[
\begin{aligned}
x_V\widehat\otimes_B y_V&=1_C,\\
y_V\widehat\otimes_C x_V&=1_B.
\end{aligned}
\tag{4.5}
\]

\textbf{Proof.} In the first product cancel
\(M_B\widehat\otimes_BM_B^*\) to the identity correspondence, and then
use \((b_Va_V)\boxtimes_R1_C=1_C\). In the second use
\((a_Vb_V)\boxtimes_R1_C=1_{B\widehat\otimes_R C\widehat\otimes_R C}\),
and cancel the dual Morita pair. These are the explicit evaluation
tensor unitaries from Lemma 4.1 and the two equations of Theorem 3.1.
Associativity is supplied by Lemma 1.1. \(\square\)

In particular, \(x_V\) has the concrete module
\(C_0(V,\pi^*\Lambda^*V_{\mathbb C})\), source Clifford action
\(i\widehat c\), grading (4.2), and position operator
\(c(v)/\sqrt{1+\|v\|^2}\). Formula (4.4) supplies its Dirac inverse.

\subsection{5. Spinors and the positive Thom
class}\label{spinors-and-the-positive-thom-class}

For an oriented Euclidean bundle, a spin\(^{c}\) structure is a lift of
its oriented orthonormal frame bundle to \[
\begin{gathered}
\operatorname{Spin}^{c}(n)\\
 =(\operatorname{Spin}(n)\times U(1))/\{\pm(1,1)\}.
\end{gathered}
\tag{5.1}
\] The spin group acts on its usual complex spin representation and
\(U(1)\) acts by scalars. The determinant line \(L\) is the associated
line for the character \([s,\lambda]\mapsto\lambda^2\).

In even rank \(n=2r\), the associated spinor bundle \(S=S^+\oplus S^-\)
has rank \(2^r\), and Clifford multiplication is odd and self-adjoint.
Its grading is the chirality \[
\Gamma_S=(-i)^r c(e_1)\cdots c(e_{2r})
\tag{5.2}
\] in a positively oriented orthonormal frame. This is the convention
established by the Pauli-matrix model in the Clifford lesson.

\textbf{Proposition 5.1 (spinor Morita equivalence).} In even rank the
spinor action identifies \[
\Gamma(\operatorname{Cl}(V))
 \cong\operatorname{End}_{C(X)}\Gamma(S).
\tag{5.3}
\] Thus \(\Gamma(S)\), and also \(\Gamma(S\otimes L^{-1})\), are graded
Morita correspondences from \(C\) to \(R\). In odd rank, a spin\(^{c}\)
structure gives a graded Morita correspondence from \(C\) to
\(R\widehat\otimes C_1\).

\textbf{Proof.} The even-rank assertion is the fibre matrix-algebra
isomorphism from the Clifford lesson, conjugated by the spinor
transition maps. It gives an isomorphism of algebra bundles and hence of
section algebras. The spinor section module is finitely generated
projective, by finite local frames and a partition of unity. Its compact
operators are all its endomorphisms. The inner-product evaluation
unitaries prove the two Morita identities as in Lemma 4.1. Twisting by a
line bundle does not change the endomorphism algebra or fullness.

For rank \(2r+1\), let \(\Delta\) be the ungraded irreducible spinor
bundle of rank \(2^r\), in the Clifford branch fixed by the orientation.
On \(\Delta\otimes C_1\), regarded as a right \(R\otimes C_1\)-module,
represent a real vector \(v\) by \(c_\Delta(v)\otimes\epsilon\), where
\(\epsilon\) is the odd generator of \(C_1\); grade only the \(C_1\)
factor. These operators are odd self-adjoint Clifford operators. Even
Clifford words give all endomorphisms of \(\Delta\) times \(1\), and odd
words give all endomorphisms times \(\epsilon\). The fibre map onto
\(\operatorname{End}(\Delta)\widehat\otimes C_1\) is therefore an
isomorphism by the equal dimensions. Spin\(^{c}\) transition maps
intertwine it. Finite generation, fullness and the evaluation identities
follow exactly as in even rank. \(\square\)

A choice of spinor convention is needed to turn (4.5) into a particular
Thom class. We choose the outward spinor symbol. In even rank it is \[
\begin{gathered}
\mathcal E_V=C_0(V,\pi^*S),\\
F_S(v)=\frac{c_S(v)}{\sqrt{1+\|v\|^2}},\\
\tau_V=[(\mathcal E_V,\pi^*,F_S)]\\
\in KK_R(R,B).
\end{gathered}
\tag{5.4}
\] Its cycle axioms follow directly from \(c_S(v)^2=\|v\|^2\), finite
rank, and compactness of the base. The following sign and line-bundle
check identifies the correct Morita reduction.

\textbf{Lemma 5.2 (spinor reduction of the Clifford Thom cycle).} Let
\(N=\Gamma(S\otimes L^{-1})\), with its \(C\)-\(R\) Morita structure.
Then \[
\tau_V=[N^*]\widehat\otimes_C x_V.
\tag{5.5}
\]

\textbf{Proof.} Work first in an oriented orthonormal fibre. The doubled
module of Lemma 4.1 decomposes under the second Clifford action as a
spinor for that action tensored with a multiplicity space. Its graded
decomposition has \(\Gamma_M=\Gamma_c\Gamma_f\), by the sorted chirality
in (4.2). After removing the source spinor, the remaining position
Clifford action therefore has grading \(\Gamma_c\). It is the positive
irreducible spinor representation, rather than its grading reverse. The
matrix-algebra classification from the Clifford lesson supplies an even
unitary identifying it with \(S\).

We check transition maps, which is necessary for a bundle statement. A
spin lift acts on the doubled module by the product of its actions in
the two Clifford copies. Indeed it conjugates both vector families by
the same orthogonal map; on exterior forms this is precisely the
exterior action. One can check this on a plane rotation: the
infinitesimal generator is \((c_jc_k-\widehat c_j\widehat c_k)/2\), the
generator of the exterior rotation. Plane rotations generate the spin
group, and the central \(-1\) acts twice, hence trivially, on the
doubled module.

The doubled module is associated to the orthogonal frame bundle, so the
scalar \(\lambda\) in (5.1) acts trivially on it. The source module
\(N=S\otimes L^{-1}\) has scalar action \(\lambda^{-1}\). Removing it
leaves scalar action \(\lambda\), exactly the action on \(S\). The
remaining spin action is the position spin representation already
identified above. Thus the fibre unitary intertwines the full
spin\(^{c}\) transition maps and is a global even bundle unitary. Under
the graded tensor evaluation it sends the position operator to
\(c_S(v)\) and the source action to pullback of \(R\). It therefore
identifies the cycles in (5.5). \(\square\)

The determinant twist in (5.5) fixes which spinor symbol is used. Using
\(S\) instead of \(S\otimes L^{-1}\) for the source Morita
correspondence would give the outward symbol twisted by \(L^{-1}\).

\textbf{Theorem 5.3 (the spin\(^{c}\) Thom isomorphism).} A spin\(^{c}\)
rank-\(n\) bundle over a compact Hausdorff base has an invertible class
\[
\tau_V\in KK_R^n(R,B).
\tag{5.6}
\] Consequently, in both degrees modulo two, \[
\begin{gathered}
K_i(C(X))\xrightarrow{\ \cong\ }K_{i+n}(C_0(V)),\\
a\longmapsto a\widehat\otimes_R\tau_V.
\end{gathered}
\tag{5.7}
\] On an oriented trivial fibre this is \((-1)^{n(n-1)/2}\) times the
ordered product of the positive one-dimensional Bott classes; in a
complex line its even symbol is \(z\) from even to odd spinors.

\textbf{Proof.} For even \(n\), combine (5.5), Proposition 5.1 and
Theorem 4.2. An explicit inverse is \[
\eta_V=y_V\widehat\otimes_C[N].
\tag{5.8}
\] The two Morita evaluation identities and (4.5) give
\(\tau_V\eta_V=1_R\) and \(\eta_V\tau_V=1_B\).

For odd rank take the ungraded spinor bundle \(\Delta\) in the
positive-volume branch, so that
\((-i)^{(n-1)/2}\prod_jc_\Delta(e_j)=1\). The right odd cycle has module
\(C_0(V,\pi^*\Delta)\widehat\otimes C_1\), graded by the last factor,
and operator \(c_\Delta(v)\otimes\epsilon/\sqrt{1+\|v\|^2}\). Its
compact defect and source commutator are exactly those of the even
position cycle. This defines the outward class \(\tau_V\).

Let \(\ell\) be a positive trivial real line. Proposition 7.2 below
proves directly from the position operators, without using
invertibility, that \(\tau_V\boxtimes_R\tau_\ell=-\tau_{V\oplus\ell}\).
Put \(u=1_B\boxtimes\tau_{\mathbb R}\) and
\(v=1_B\boxtimes\eta_{\mathbb R}\), with \(uv=1_B\) and
\(vu=1_{C_0(V\oplus\ell)}\). The even case gives the inverse
\(\eta_{V\oplus\ell}\). Consequently \(\tau_V=-\tau_{V\oplus\ell}v\) and
\(\eta_V=-u\eta_{V\oplus\ell}\) are inverse classes. Associating these
two displayed products proves both inverse identities. This is the
ordered addition and removal of a last line, including its required
sign.

Here is the orientation check in those maps. For one real line the
doubled model has source \(i\widehat c=\sigma_2\), operator
\(h(v)\sigma_1\), and grading \(\sigma_3\), where
\(h(v)=v/\sqrt{1+v^2}\). The even unitary \(\exp(i\pi\sigma_3/4)\) sends
it to source \(\sigma_1\) and operator \(-h(v)\sigma_2\). Proposition
7.1 of the Bott lesson identifies precisely this left-Clifford cycle
with the positive cone class. The right odd transfer therefore gives the
positive one-dimensional class. When several odd classes are multiplied,
the course's right-Clifford product orders the second cycle's original
generator before the first cycle's auxiliary generator. Proposition 7.2
below computes the resulting sign: on a positively oriented
\(\mathbb R^n\), the outward spinor class is \((-1)^{n(n-1)/2}\) times
the ordered exterior product of the positive one-dimensional classes.
Adding and removing a last line uses this same sign. The graded Clifford
equivalence before scalar transfer has no such scalar-degree conversion.

In a complex line the spinor model is \(S^+=\mathbb C\),
\(S^-=\mathbb C\), with \[
c_S(x+iy)=
\begin{pmatrix}0&x-iy\\x+iy&0\end{pmatrix}.
\tag{5.9}
\] Its lower-left symbol is \(z=x+iy\). The explicit self-adjoint disc
lift in
\href{supporting/relative-k-foundations/relative-k-foundations.html\#rk-planar-bott}{Theorem
RK.6} therefore gives exactly the positive planar class
\([q]-[\operatorname{diag}(0,1)]\) in that comparison convention. Graded
tensor products of these models prove the complex higher-rank assertion.

Finally multiply K-theory classes by the inverse pair. The associativity
proof applies by Lemma 1.1; for scalar-source cycles one can adjoin the
central right \(R\)-action on the essential scalar summand. Clifford
dilation gives the same argument in odd degree. This proves bijectivity
in (5.7). For a vector bundle \(E\) representing \(a\in K^0(X)\), the
tensor module is \(\pi^*E\otimes\pi^*S\) and its operator is
\(1_E\otimes c_S(v)/\sqrt{1+\|v\|^2}\). Thus the formula is the usual
\(\pi^*a\) times the Thom symbol, including bundle differences.
\(\square\)

\subsection{6. Complex bundles and relative
symbols}\label{complex-bundles-and-relative-symbols}

\textbf{Proposition 6.1 (the complex spinor bundle).} A Hermitian
complex rank-\(r\) bundle \(E\), viewed as a real rank-\(2r\) bundle,
has a spin\(^{c}\) structure with \[
\begin{gathered}
S=\Lambda^*E,\quad L=\det E,\\
c(v)=\varepsilon_v+\iota_v.
\end{gathered}
\tag{6.1}
\] Its Thom class is the relative triple \[
\begin{gathered}
\left(\pi^*\Lambda^{\rm even}E,\pi^*\Lambda^{\rm odd}E,\sigma\right),\\
\sigma=c(v)^+\big|_{S(E)}.
\end{gathered}
\tag{6.2}
\] The comparison on the sphere is unitary. At the zero section \(s\),
\[
\begin{gathered}
s^*\tau_E=\lambda_{-1}(E)\\
=\sum_{j=0}^r(-1)^j[\Lambda^jE].
\end{gathered}
\tag{6.3}
\]

\textbf{Proof.} Exterior multiplication and its Hermitian adjoint give
\(\{c(v),c(w)\}=2\operatorname{Re}\langle v,w\rangle\). In a unitary
frame with ordered real basis \(e_1,ie_1,\ldots,e_r,ie_r\), these are
the tensor Pauli generators. They generate every endomorphism of
\(\Lambda^*E\), and their chirality is its form parity. Hence this is an
irreducible graded Clifford module.

For completeness, this supplies an actual spin\(^{c}\) lift. The group
of pairs \((g,u)\), with \(g\in SO(2r)\) and \(u\) an even unitary
implementing \(c(gv)\), is \(\operatorname{Spin}^{c}(2r)\). A spin lift
supplies one such unitary for every \(g\); two implementers differ by a
scalar because the Clifford matrices generate the full endomorphism
algebra. The map from \(\operatorname{Spin}(2r)\times U(1)\)
consequently has precisely the kernel \(\{(1,1),(-1,-1)\}\). The unitary
group \(U(r)\) maps continuously to this implementer group by
\(g\mapsto(g,\Lambda^*g)\). Applying it to the unitary frame transitions
gives the required lift.

Its determinant character is \(\det g\). On a diagonal unitary rotation
of angles \(\theta_j\), the exterior representation differs from the
spin representation by the scalar \(\exp(i\sum_j\theta_j/2)\), whose
square is \(\det g\). Every unitary is unitarily diagonalizable;
conjugation preserves both characters, proving the equality for all
\(g\). This proves \(L=\det E\).

The compact Hausdorff ball/sphere identification and both inverse
relative-bundle maps in
\href{supporting/relative-k-foundations/relative-k-foundations.html\#rk-difference}{Theorems
RK.3--RK.4} turn the outward symbol (5.4) into (6.2).
\href{supporting/relative-k-foundations/relative-k-foundations.html\#rk-symbol-rules}{Corollary
RK.4a} proves the symbol and zero-section rules for these actual
comparisons. On the unit sphere \(c(v)^2=1\), so its even-to-odd block
is unitary. Restriction to the zero section forgets that boundary
comparison and gives the difference of the two bundles, namely (6.3).
\(\square\)

For a line bundle the triple is simply \[
(\mathbf1,\pi^*E,\ \lambda\mapsto\lambda v).
\tag{6.4}
\] This formula keeps track of the actual bundle and comparison, rather
than just the fibre winding.

\subsection{7. Pullback and direct sums}\label{pullback-and-direct-sums}

\textbf{Theorem 7.1 (naturality).} Let \(f:Y\to X\) be a continuous map
of compact Hausdorff spaces. Give \(f^*V\) the pulled-back metric and
spin\(^{c}\) structure. Its total-space map \(\widetilde f:f^*V\to V\)
is proper, and \[
\widetilde f^*(\tau_V)=\tau_{f^*V}
\tag{7.1}
\] as Thom symbols. In K-theory, \[
\widetilde f^*\bigl(\pi^*a\,\tau_V\bigr)
 =\pi_Y^*f^*a\,\tau_{f^*V}.
\tag{7.2}
\] The Clifford equivalence and its inverse pull back as well.

\textbf{Proof.} A compact subset of \(V\) is norm bounded. Its inverse
image under \(\widetilde f\) is closed in a norm-bounded ball bundle
over the compact base \(Y\), hence compact. This proves properness and
gives the pullback homomorphism on \(C_0\)-algebras.

Pulling back the position module gives the module of sections of the
pulled-back spinor bundle, and its operator is still
\(c(v)/\sqrt{1+\|v\|^2}\). The fibre isometry in the definition of
\(f^*V\) preserves \(v\), its norm, and both Clifford actions. Thus the
cycle identification proves (7.1), including in the odd model with the
additional last line.

More explicitly, the position module scalar-extension map sends
\(\xi\otimes b\) to \((\xi\circ\widetilde f)b\). Its inner-product
formula is that for scalar extension. Its range is dense by local frames
and finite partitions, so it is the claimed module unitary. The same
argument for the finite Morita modules preserves (4.2).

For the inverse, scalar extension of the vertical \(L^2\) module has the
pulled-back Hilbert field. Local smooth fibre sections are dense in both
modules; on them the fibre-integral inner products, Dirac operators and
inverse resolvents agree. This proves the inverse-cycle identification
and its graph-domain compatibility. Every homotopy in Theorem 3.1 has
these same frame-invariant formulas, so the inverse identities also pull
back.

For a bundle representative of \(a\), tensor its pullback with the
pulled-back spinor symbol. This gives (7.2) directly; differences and
Clifford suspension give both degrees. Equivalently, the Clifford
statement is the KK equality \[
[C_f]\widehat\otimes_{C_{f^*V}}x_{f^*V}
 =x_V\widehat\otimes_B[B_f],
\tag{7.3}
\] where \(C_f\) and \(B_f\) are the Clifford and total-space pullback
homomorphisms. Both sides have the same position cycle just identified.
\(\square\)

Metric choices cause no ambiguity. Join two metrics through their convex
interpolation, and transport Clifford vectors by the positive fibrewise
isometry between them. The same formulas give a cycle on the compact
base cylinder. Endpoint evaluation proves metric independence with the
given spin\(^{c}\) structure.

\textbf{Proposition 7.2 (ordered direct sums).} For spin\(^{c}\) bundles
\(V,W\) of ranks \(r,s\) over \(X\), using the ordinary oriented
direct-sum spinor convention with \(V\) first and the course's
right-Clifford degree products, \[
\tau_{V\oplus W}=(-1)^{rs}\tau_V\boxtimes_R\tau_W.
\tag{7.4}
\] For a trivial bundle this is \(1_X\boxtimes\tau_{\mathbb R^n}\).

\textbf{Proof.} In even rank the spinor module is the graded tensor
product, and the unbounded symbol is \[
c_V(v)\widehat\otimes1+1\widehat\otimes c_W(w).
\tag{7.5}
\] Its square is \(\|v\|^2+\|w\|^2\), since the two odd summands
anticommute. Its explicit inverse resolvents prove regularity. The
domain is contained in each summand's domain. Creations from compactly
supported smooth first-factor spinors have bounded error, namely
creation by \(c_V(v)\xi\), and the positivity form for the first summand
is \(2\|v\|^2\geq0\). Its inverse resolvent is locally compact on the
finite-rank symbol module. The product criterion identifies its bounded
transform with the external product. It is exactly the direct-sum
symbol. The central argument applies over any compact \(X\).

We spell out the odd transfer. If just one rank is odd, the scalar odd
operator is the usual direct-sum spinor position operator: for an even
first factor it is \(c_V(v)\otimes1+\Gamma_V\otimes c_W(w)\); for an
even second factor it is \(c_V(v)\otimes\Gamma_W+1\otimes c_W(w)\).
Ordered volume multiplication gives the positive odd spinor convention.
Their squares, creation errors and first-factor positivity are the same
ones checked above, in the right odd Clifford module. There is no sign
in these two cases.

If both ranks are odd, use the ungraded positive-volume spinor bundles
\(S_V,S_W\). The ordinary oriented direct sum has bundle
\(S_V\otimes S_W\otimes\mathbb C^2\), grading \(\sigma_3\), and position
operator \[
Q_+=c_V(v)\otimes\sigma_1+c_W(w)\otimes\sigma_2.
\tag{7.6}
\] The omitted tensor identities act on the other spinor factor. Write
\(r=2a+1,s=2b+1\). The positive odd volumes satisfy \(\prod c_V=i^a\),
\(\prod c_W=i^b\); multiplying the ordered total volume gives
\((-i)^{a+b+1}i^{a+b}\sigma_1\sigma_2=\sigma_3\). This verifies the
stated direct-sum grading globally. Spin transitions act on their own
two spinor factors, so this calculation is a bundle identification.

In the product, however, the original second Clifford generator comes
first and the auxiliary first generator second. Their matrices are
\(-\sigma_2,\sigma_1\), exactly as in Proposition 8.1 of
\emph{Connections and the existence of the Kasparov product}.
Consequently its position operator is \[
\begin{gathered}
Q_-=c_V(v)\otimes\sigma_1-c_W(w)\otimes\sigma_2,\\
\Gamma=\sigma_3.
\end{gathered}
\tag{7.7}
\] Its square is again \(\|v\|^2+\|w\|^2\). Dense compactly supported
creating sections have the fixed first position error; the Clifford
anticommutation gives their creation sign and the adjoint error. The
first-operator domain inclusion and positivity form \(2\|v\|^2\) hold as
before. The product criterion therefore gives this actual cycle.
Conjugation by \(1\otimes\sigma_1\) sends \(Q_-\) to \(Q_+\) and sends
its grading to \(-\sigma_3\). Thus its class is the additive inverse of
the ordinary direct-sum Thom class. This proves the minus sign when both
ranks are odd and hence (7.4) in every parity.

For two positive lines this is the elementary distinction between the
symbols \(x-iy\) and \(x+iy\). Iteration gives the factor
\((-1)^{n(n-1)/2}\) stated in Theorem 5.3. For a trivial bundle all
modules are the corresponding exterior scalar extensions of the actual
fibre module, proving the last assertion. \(\square\)

\subsubsection{Corollary 7.3. Translating an ordered-coordinate
normalization}\label{corollary-7.3.-translating-an-ordered-coordinate-normalization}

Set \(\epsilon_n=(-1)^{n(n-1)/2}\). In the right-Clifford product
convention of this course, the rank-normalized symbols \[
\widehat\tau_V=\epsilon_r\tau_V,
\qquad r=\operatorname{rank}V,
\] satisfy \[
\widehat\tau_{V\oplus W}
=\widehat\tau_V\boxtimes_R\widehat\tau_W.
\tag{7.8}
\] For an oriented trivial bundle, \(\widehat\tau_{\mathbb R^n}\) is the
ordered right-Clifford external product of its positive one-dimensional
symbols. Equivalently, one may retain the outward symbols \(\tau_V\) and
use the degree product \[
\begin{gathered}
x\boxtimes_{\mathrm{out}}y
=(-1)^{rs}x\boxtimes_R y,\\
|x|=r,\qquad |y|=s.
\end{gathered}
\tag{7.9}
\] Then the outward direct-sum symbols multiply without an additional
sign.

\textbf{Proof.} The integer identity \[
\begin{aligned}
\frac{(r+s)(r+s-1)}2
&=\frac{r(r-1)}2\\
&\quad+\frac{s(s-1)}2+rs.
\end{aligned}
\] gives \(\epsilon_{r+s}=\epsilon_r\epsilon_s(-1)^{rs}\). Substitute
Proposition 7.2 into \(\epsilon_{r+s}\tau_{V\oplus W}\); the two factors
\((-1)^{rs}\) cancel. This proves (7.8). Starting with \(\epsilon_1=1\)
and applying it successively to ordered trivial lines proves the
coordinate assertion. Formula (7.9) and Proposition 7.2 give the second
formulation directly. The product in (7.9) is associative: the signs for
three degrees on its two sides are \((-1)^{rs+(r+s)t}\) and
\((-1)^{st+r(s+t)}\), which agree. It has the same naturality and
bilinearity as the original product.

The degreewise conversion is explicit: multiplication by \(\epsilon_n\)
takes the right-Clifford product to (7.9), since \[
\epsilon_{r+s}(x\boxtimes_R y)
=(\epsilon_r x)\boxtimes_{\mathrm{out}}(\epsilon_s y).
\] One must retain the integer degree while making this conversion.
Although KK is two-periodic, \(\epsilon_{n+2}=-\epsilon_n\); changing
the convention therefore changes the chosen degree-two periodicity
identification as well. A sign cannot be discarded merely by replacing a
rank with its parity. \(\square\)

In rank two the distinction is visible before any Chern-character
calculation. Two positive lines give the right-product operator
\(x\sigma_1-y\sigma_2\), while the outward plane has
\(x\sigma_1+y\sigma_2\), with grading \(\sigma_3\). Proposition 7.2
proves their classes are opposite by an explicit grading-reversing
conjugation. Thus \(\epsilon_2=-1\). The next values are
\(\epsilon_3=-1\), \(\epsilon_4=1\). These comparisons specify the
actual product and periodicity conventions, rather than using the word
``positive'' alone to identify them.

\subsection{8. The Chern character and the Todd
factor}\label{the-chern-character-and-the-todd-factor}

The following comparison uses rational cohomology on a finite CW base,
or on a compact base of finite CW homotopy type. The KK Thom theorem
above requires neither assumption. Let \(U_E\) be the ordinary
cohomological Thom class for the complex orientation, normalized to
integrate to \(+1\) on each real \(2r\)-dimensional fibre. We identify
its group \(H^*(E,E\setminus s(X))\) with \(H^*(D(E),S(E))\): the disc
bundle is homotopy equivalent to \(E\), and the sphere bundle is
homotopy equivalent to the punctured bundle, so their pair sequences
give the identification.

We use the complete local proofs in
\href{supporting/cohomological-foundations/cohomological-foundations.html}{Cohomological
foundations for the complex Thom comparison}:
\href{supporting/cohomological-foundations/cohomological-foundations.html\#cf-thom}{CF.1b--CF.1c}
prove the ordinary Thom isomorphism, its naturality and its ordered
product;
\href{supporting/cohomological-foundations/cohomological-foundations.html\#cf-splitting}{CF.4a}
proves the splitting pullback injection; and
\href{supporting/cohomological-foundations/cohomological-foundations.html\#cf-character}{CF.5a--CF.5b}
construct the even relative Chern character and prove its products and
bundle action. Its normalization is \(\operatorname{ch}(L)=e^{c_1(L)}\);
\href{supporting/cohomological-foundations/cohomological-foundations.html\#cf-projective-space}{CF.3a
and CF.5b} prove that the positive planar Bott projection has Chern
value \(-1\). Only this natural even character is needed; no rational
comparison isomorphism on arbitrary compact Hausdorff spaces is used.

Define the Todd class by its line factors \[
\operatorname{Td}(E)=\prod_j\frac{x_j}{1-e^{-x_j}},
\tag{8.1}
\] where \(x_j\) are the Chern roots after the splitting pullback. The
factors are power series with constant one; positive-degree classes are
nilpotent, so only finitely many terms contribute.

\textbf{Proposition 8.1 (the factor in the fixed convention).} For the
complex spinor symbol (6.2), \[
\begin{gathered}
\operatorname{ch}_c(\tau_E)\\
 =(-1)^r\,\pi^*g_E\smile U_E,\\
g_E=e^{c_1(E)}\operatorname{Td}(E)^{-1}.
\end{gathered}
\tag{8.2}
\]

\textbf{Proof.} First take a line \(L\) with \(x=c_1(L)\). The
cohomological Thom isomorphism writes its compactly supported character
uniquely as \(\pi^*g(x)U_L\), with an even-degree coefficient on the
base. By (6.3), pulling back to the zero section gives \[
x\,g(x)=1-e^x.
\tag{8.3}
\] We justify determination of the coefficient instead of dividing a
possibly zero-divisor in a particular base. On the tautological line
over \(\mathbb{CP}^N\), the ring calculation in
\href{supporting/cohomological-foundations/cohomological-foundations.html\#cf-projective-space}{CF.3a}
makes all coefficients polynomials in \(x=-h\). Equation (8.3)
determines their terms below the top degree. Pull back the identical
construction over \(\mathbb{CP}^{N+1}\); its next relation determines
that last term as well. Naturality of the relative Chern character in
CF.5b and of the ordinary Thom class in CF.1c makes these coefficients
compatible. A line bundle over a finite CW base is pulled back from a
finite projective space by the finite isometric embedding proved in
\href{supporting/relative-k-foundations/relative-k-foundations.html\#rk-line-classification}{Lemma
RK.7a}. Thus all line bundles have the coefficient \[
g(x)=\frac{1-e^x}{x}
 =-e^x\frac{1-e^{-x}}{x}.
\tag{8.4}
\] The expression at \(x=0\) means its power-series value \(-1\). This
also agrees with the actual planar projection's Chern number.

The relative Chern character preserves these actual symbol products by
\href{supporting/cohomological-foundations/cohomological-foundations.html\#cf-relative-character}{CF.5b}.
Its proof constructs the graded tensor comparison, verifies the
neighborhood deformation needed at the quotient wedge, and proves
injectivity of the smash-quotient pullback from the pair sequence.
Ordinary bundle-character multiplicativity then gives the relative
product, and base-diagonal pullback gives the balanced symbol product.
These are the disc/sphere pairs covered by that proof.

Now apply the splitting pullback proved in
\href{supporting/cohomological-foundations/cohomological-foundations.html\#cf-splitting}{CF.4a}.
It is injective on base cohomology; the two ordinary Thom isomorphisms
imply that it is injective on the corresponding relative bundle
cohomology too. On this base \(E\) is a sum of lines. Proposition 7.2
multiplies their K-symbols, and the ordinary Thom classes multiply in
the same ordered complex orientation. Multiplying (8.4) gives \[
\prod_j\frac{1-e^{x_j}}{x_j}
 =(-1)^r e^{\sum_jx_j}\operatorname{Td}(E)^{-1}.
\] Injectivity brings this equality back to the original base and proves
(8.2). \(\square\)

If one instead uses the dual Koszul symbol with zero-section class
\(\lambda_{-1}(E^*)\), the same line calculation gives \((1-e^{-x})/x\).
Multiplication after the splitting pullback gives its factor
\(\operatorname{Td}(E)^{-1}\), with fibre Chern value \(+1\). Formula
(8.2) records both changes for our outward symbol: its zero-section
class is \(\lambda_{-1}(E)\), and its ordered positive fibre class has
Chern value \((-1)^r\).

\subsection{9. Compact group actions}\label{compact-group-actions}

We supply the compact equivariant operations needed here before using
them. Thus the inverse identities in this section depend only on
preceding programme proofs and the following averaging argument.

For a second-countable compact group \(G\), a graded equivariant Hilbert
\(B\)-module has an even strongly continuous action \(U_g\) with
\(U_g(\xi b)=U_g\xi\,\beta_g(b)\) and
\(\langle U_g\xi,U_g\eta\rangle=\beta_g\langle\xi,\eta\rangle\). An
equivariant cycle is an ordinary graded Kasparov cycle with equivariant
representation \(\phi\), norm-continuous orbit
\(g\mapsto U_gFU_g^{-1}\), and extra localized defect
\((U_gFU_g^{-1}-F)\phi(a)\in\mathcal K(E)\). Equivariant homotopies are
cycles over the interval coefficient algebra, with trivial action on the
interval. Their homotopy classes, with direct sum, define \(KK^G\). The
opposite-grading and rotation proofs of the preceding cycle lesson give
additive inverses: their scalar matrices commute with the group action,
and a degenerate cycle contracts by the C\_0-path construction of
Proposition 2.3 of Lesson 06, with diagonal action. Auxiliary Clifford
factors have trivial action.

\textbf{Lemma 9.0 (compact equivariant products).} For separable graded
G-algebras and countably generated equivariant modules, the connection
and localized positivity criterion of \emph{Connections and the
existence of the Kasparov product} defines a bilinear equivariant
product. It is associative and has equivariant homomorphism cycles as
identities. The exterior extensions and signed flips of \emph{Homotopy,
associativity, the index pairing and KK-equivalence} are equivariant. An
invariant unbounded candidate satisfying the local Theorem 0.5
represents this product.

\textbf{Proof.} Write \(g\cdot T=U_gTU_g^{-1}\). Rank-one operators
satisfy \(g\cdot\theta_{\xi,\eta}=\theta_{U_g\xi,U_g\eta}\), so compact
operators have norm-continuous orbits. The orbit of an adjointable
operator is continuous on every vector, as is its adjoint orbit: this
follows by inserting \(U_g^{-1}\xi\) and using strong continuity and the
uniform operator bound. A bounded such field has a strict integral,
obtained by integrating its action and its adjoint action on vectors;
these actions respect the right module and the adjoint identity. Haar
translation makes its average invariant.

First make an input operator self-adjoint and contractive by the
preceding cycle normalization, whose scalar functional calculus
preserves the equivariance defect. Its Haar average \(\bar F\) is
invariant, self-adjoint and contractive. For each \(a\), \[
(\bar F-F)\phi(a)=\int_G(g\cdot F-F)\phi(a)\,dg
\] is compact: the integrand is a norm-continuous compact field. The
adjoint formula supplies \(\phi(a)(\bar F-F)\in\mathcal K(E)\).
Expanding the commutator and square shows that \((1-t)F+t\bar F\) is an
equivariant cycle for every \(t\). Its localized equivariance error is
\((1-t)(g\cdot F-F)\). Its orbit is norm continuous. This also works on
interval modules, so classes and homotopies can be represented
invariantly.

For two invariant inputs, choose an ordinary self-adjoint odd connection
\(G_0\) by Theorem 1.3 of the product lesson and average it strictly.
For homogeneous \(\xi\), conjugation sends the creation operator
\(T_\xi\) to \(T_{U_g\xi}\). The creation error of \(g\cdot G_0\) is
therefore the conjugate of the error of \(G_0\) for \(U_g^{-1}\xi\). It
varies norm continuously with \(g\): creation operators are norm
continuous in their vectors, and conjugation is norm continuous on
compact operators. Both creation errors have compact integrals. The
averaged \(G\) is consequently an invariant self-adjoint odd connection.

The only further change to the ordinary existence proof is an invariant
technical partition. Suppose its separation data \(J,A_1,A_2,\Delta\)
are invariant and consist of norm-continuous orbit elements. Apply
Theorem 3.1 of \emph{Kasparov's technical theorem}, then strictly
average its even positive contractions \(M_0,N_0\). For \(a\in A_1\) and
\(d\in\Delta\), the identities \[
(g\cdot M_0)a=g\cdot\bigl(M_0(g^{-1}\cdot a)\bigr),\qquad
[g\cdot M_0,d]=g\cdot[M_0,g^{-1}\cdot d]
\] give norm-continuous \(J\)-valued fields. Their integrals prove
\(MA_1\subset J\), \([M,\Delta]\subset J\); the identical calculation
proves \(NA_2\subset J\). Positivity, parity, \(M+N=1\), and invariance
are preserved. Corollary 3.2 of that lesson supplies the required root
properties by functional calculus in the quotient.

Use exactly the separation data in Theorem 3.1 of the product lesson:
with \(S=F_1\widehat\otimes1\), they are \[
\begin{aligned}
J&=\mathcal K(E_1\widehat\otimes_D E_2),\qquad
I=\mathcal K(E_1)\widehat\otimes1,\\
A_1&=C^*(J,I),\\
A_2&=C^*(J,G^2-1,[S,G]_{\mathrm{gr}},
 [G,\phi(A)]_{\mathrm{gr}}),\\
\Delta&=\overline{\operatorname{span}}
 \{S,G,\phi(A),S^*,G^*,\phi(A)^*\}.
\end{aligned}
\] Invariance of \(S,G\), equivariance of \(\phi\) and invariance of the
compact images make these data invariant. Their orbit continuity follows
from their displayed generators. The countability, separation and
derivation conditions are proved in that theorem before its partition
step; none uses a group action. The averaged partition therefore gives
the invariant candidate \(F=M^{1/4}SM^{1/4}+N^{1/4}GN^{1/4}\). Equations
(3.5)--(3.7) of that proof verify its cycle defects, connection errors
and localized positivity. Its additional equivariance defect is zero.

For invariant product candidates \(F,F'\), use the comparison algebra
\(C^*(J,[S,F]_{\mathrm{gr}},[S,F']_{\mathrm{gr}},F-F')\) and include
\(S,F,F',\phi(A)\) and their adjoints in the derivating space. These are
invariant separation data by the zero-connection calculation in Theorem
4.1 of the product lesson. The averaged partition gives the invariant
comparison operator of its equation (4.3); its positive-anticommutator
normalization paths are invariant functional calculus. This proves
invariant uniqueness. A possibly noninvariant equivariant product can
first be averaged: its connection errors integrate as above, its
localized positivity integrates to a positive element of the
compact-operator quotient, and the two-sided locally compact straight
homotopy just proved identifies its class with the average. Thus
uniqueness holds among all equivariant products.

Construct a product over an interval module to prove homotopy invariance
in either variable; evaluation gives the endpoint products by their same
creation and positivity equations. Block sums prove bilinearity. For an
equivariant homomorphism, the essential-replacement interval ideal of
Proposition 5.1 of the product lesson is invariant. Choose its
connection by averaging. Its endpoint tensor maps are equivariant, so
Theorem 5.2 of that lesson gives both identity laws even for degenerate
representations.

For three inputs use the proof of Theorem 4.1 of the
homotopy-and-associativity lesson, with invariant intermediate product
operators. The compact images, bracket connections and negative parts in
that proof are invariant. Replace its technical partition by its Haar
average. The two terms of its positivity equation (4.3) remain positive:
one uses the first intermediate product's positive lift and annihilation
of its compact error, and the other uses annihilation of the final
bracket's negative part. The simultaneous candidate is now invariant, so
uniqueness proves associativity. The exterior tensor and flip unitaries
are equivariant for diagonal actions; the same uniqueness comparison
proves the signed exterior laws. Finally the resolvent connection and
integrated positivity estimates of the unbounded criterion are equations
on these same modules. For an invariant candidate its resolvents,
bounded transform and normalization homotopies are invariant; the
estimates verify the two product conditions just proved sufficient. This
proves the last assertion. \(\square\)

\textbf{Theorem 9.1 (equivariant linear Thom theorem).} Let a compact
group \(G\) act continuously and linearly on a finite-dimensional real
vector space \(W\). Then there are inverse classes \[
\begin{gathered}
x_W\in KK^G(\operatorname{Cl}(W),C_0(W)),\\
y_W\in KK^G(C_0(W),\operatorname{Cl}(W)).
\end{gathered}
\tag{9.1}
\] If the action has a spin\(^{c}\) lift, it gives a degree-\(\dim W\)
equivariant Thom equivalence with the scalar algebra. In particular
\(C_0(W)\) with this action is equivariantly KK-equivalent to \(C_0(W)\)
with the trivial action. A complex linear action has the required lift.

\textbf{Proof.} Average an inner product over normalized Haar measure,
making the action orthogonal. Give
\(L^2(W)\otimes\Lambda W_{\mathbb C}\) its change-of-variable and
exterior representation. Position Clifford multiplication and the
vertical Dirac operator intertwine this action exactly. Their inverse
resolvents and bounded transforms are invariant. Thus the cycles of
Sections 2--4 are equivariant, with zero equivariance defect.

The oscillator product is equivariant. The Gaussian kernel is the
trivial one-dimensional representation, because its norm is radial and
its form degree is zero. The complementary sign homotopy is invariant
functional calculus. Hence the first product is the equivariant
identity. The simultaneous coordinate and Clifford rotation commutes
with the diagonal orthogonal action on \(W\oplus W\), at every
parameter. Its uniform compactness estimate is unchanged. The
invertible-idempotent argument consequently proves the reverse
equivariant product. The doubled Clifford Morita correspondence has the
exterior action and the invariant grading (4.2), so its two evaluation
tensor maps are equivariant. This proves (9.1).

A spin\(^{c}\) lift gives the equivariant spinor and determinant-line
correspondences of Section 5; every transition-map computation there is
an intertwining computation. They supply the degree-shifted scalar
equivalence. The trivial linear action has the same degree-shifted
scalar equivalence, so composing one with the inverse of the other gives
the asserted equivalence between the two function algebras. For a
complex linear action average a Hermitian metric and use the
homomorphism \(U(r)\to\operatorname{Spin}^{c}(2r)\) constructed in
Proposition 6.1.

There is no hidden second-countability restriction on the compact group
in this linear assertion. Its action factors through a closed subgroup
of the finite-dimensional orthogonal group, hence through a compact
metrizable group. In the spinor case include its finite-dimensional
spinor representation in that quotient. The equivariant products just
used can be taken in this second-countable quotient, where Lemma 9.0
applies, and pulled back along \(G\)'s quotient homomorphism. Pullback
leaves each creation equation, positivity sandwich and homotopy
unchanged and supplies a continuous action by composition. Their
explicit cycles and homotopies are still equivariant for \(G\), and
their inverse identities remain the same identities. \(\square\)

\subsection{10. Examples}\label{examples}

For \(X\times\mathbb C^r\), take \(S=X\times\Lambda^*\mathbb C^r\). The
symbol is the graded sum of the coordinate symbols (5.9), so Theorem 5.3
is external product with the fixed positive Bott class. In degree zero
its inverse sends that class back to the trivial line on \(X\).

Let \(E_k\to S^2\) be a complex line with \(c_1(E_k)=k h\), where \(h\)
evaluates to one on the complex-oriented sphere. Its Thom triple is
\((\mathbf1,\pi^*E_k,v)\). Its zero-section class is \(1-[E_k]\); the
comparison map contains the tautological fibre vector, so is invertible
off the zero section. Formula (8.2) gives \[
\begin{gathered}
\operatorname{ch}_c(\tau_{E_k})\\
 =-\left(1+\frac{k}{2}\pi^*h\right)U_{E_k}.
\end{gathered}
\tag{10.1}
\] Here \(h^2=0\). The Thom group is free of rank two, because
\(K^0(S^2)\cong\mathbb Z^2\), and its odd group is zero. The symbol
generates it as a \(K^0(S^2)\)-module, which is stronger than saying its
restriction generates each fibre group.

For a smooth embedding \(i:M\hookrightarrow N\) with compact \(M\), a
spin\(^{c}\) normal bundle \(\nu\) gives the first part of the wrong-way
construction: \[
K^*(M)\xrightarrow{\ \tau_\nu\ }
K_c^{*+\operatorname{rank}\nu}(\nu).
\tag{10.2}
\] A tubular neighborhood identifies this compact-support group with
that of the neighborhood in \(N\), and extension by zero gives a class
in \(K^{*+\operatorname{rank}\nu}(N)\). The composition and independence
questions for general wrong-way maps are proved in the lesson
\emph{Wrong-way maps for K-oriented maps}. The Thom step (10.2),
including its inverse and orientation, is established here.

\subsection{11. Exercises}\label{exercises}

\textbf{11.1. A trivial bundle.} Write the Clifford Thom cycle and the
spinor Thom cycle for \(X\times\mathbb R^n\). Prove that they are the
external scalar extensions of the corresponding fibre cycles.

\textbf{11.2. Complex spinors.} For a Hermitian rank-\(r\) complex
bundle \(E\), construct its Clifford spinor module, prove the Morita
equivalence, and compute its determinant line and zero-section Thom
class.

\textbf{11.3. Pullback.} Prove the naturality formula (7.2), including
properness of the total-space pullback map. Explain why the same proof
identifies the inverse Dirac cycle.

\textbf{11.4. The tautological line.} Let \(L\to\mathbb{CP}^1\) be the
tautological line and put \(u=1-[L]\). Give its Thom class as an actual
relative difference-bundle triple, compute the two compactly supported
K-groups of its total space with generators, and compute zero-section
restriction on those generators. Compare the fibre sign with the
relative-bundle convention.

\subsection{12. Solutions}\label{solutions}

\textbf{Solution to 11.1.} All frame transitions are identity. The
Clifford cycle has module
\(C_0(X\times\mathbb R^n,\Lambda^*\mathbb C^n)\), source
\(C(X)\widehat\otimes C_n\) acting by the scalar base function and
\(i\widehat c\), operator \(c(v)/\sqrt{1+\|v\|^2}\), and grading (4.2).
These are precisely the fibre cycle's scalar extensions. Its inverse is
the module \(C(X,L^2(\mathbb R^n)\otimes\Lambda^*\mathbb C^n)\),
followed by the finite Clifford Morita map in (4.4); its inverse
resolvents are the constant fibre resolvents. With the trivial positive
spinor bundle, the spinor cycle is \(C_0(X\times\mathbb R^n,S_n)\) with
outward Clifford symbol. The scalar-extension tensor maps preserve inner
products and have dense range by finite sums of base functions times
fibre sections. Hence they are module unitaries identifying the cycles.
Proposition 7.2 and its ordered one-dimensional check give
\(1_X\boxtimes\tau_{\mathbb R^n}\).

\textbf{Solution to 11.2.} Set \(S=\Lambda^*E\), grade by degree, and
use \(c(v)=\varepsilon_v+\varepsilon_v^*\). The exterior relations give
its self-adjoint Clifford relations for the underlying real Euclidean
metric. In a unitary frame, the real pair \(e_j,ie_j\) gives
\(\varepsilon_j+\iota_j\) and \(i(\varepsilon_j-\iota_j)\). Recovering
\(\varepsilon_j,\iota_j\) and using the vacuum projection gives every
matrix unit, so
\(\operatorname{Cl}(E_{\mathbb R})\cong\operatorname{End}(S)\). Finite
frames and partitions make \(\Gamma(S)\) finitely generated projective
and full; its dual and inner-product evaluation give the inverse Morita
correspondence. The implementer lift is \(g\mapsto(g,\Lambda^*g)\). Its
scalar-square character is \(\det g\), as checked on diagonal unitary
rotations in Proposition 6.1, so the determinant line is \(\det E\). The
sphere comparison is \(c(v)^+\), and the zero-section difference is
\(\sum_j(-1)^j[\Lambda^jE]\).

\textbf{Solution to 11.3.} The pulled-back metric has
\(\|v\|=\|\widetilde f(v)\|\). The inverse image of a compact set is
closed and bounded in a ball bundle over compact \(Y\), so is compact.
Pullback therefore acts on the function algebras. The pulled-back spinor
module and position symbol are exactly those of \(f^*V\); local frames
prove the scalar-extension map is an onto isometry. Pulling back a
bundle representative of \(a\) identifies the tensor symbol with
\(\pi_Y^*f^*a\) times this same symbol, proving (7.2). For the Dirac
inverse, the fibre integral, smooth fibre core, Clifford multiplication
and inverse resolvents agree on pulled-back local sections. Those
sections are dense, so the module unitary also intertwines the closures
and bounded transforms. This proves the inverse identification, rather
than inferring it only from a fibrewise K-group calculation.

\textbf{Solution to 11.4.} On the closed ball and sphere bundle the
triple is \[
\xi=(\mathbf1,\pi^*L,\
\sigma_{(\ell,v)}:\lambda\mapsto\lambda v).
\tag{12.1}
\] For \(\|v\|=1\) it is a unitary isomorphism from the trivial line to
the pulled-back tautological line. This is the specified comparison in
the difference-bundle theorem; the bundles alone would not specify the
compact-support class.

The integral clutching calculation in
\href{supporting/relative-k-foundations/relative-k-foundations.html\#rk-sphere-ring}{Theorem
RK.7} gives \[
\begin{gathered}
K^0(\mathbb{CP}^1)=\mathbb Z[1]\oplus\mathbb Z u,\\
K^1(\mathbb{CP}^1)=0,\quad u^2=0.
\end{gathered}
\tag{12.2}
\] For the last relation, the rotation path in
\href{supporting/relative-k-foundations/relative-k-foundations.html\#rk-sphere-ring}{Theorem
RK.7} joins the clutching matrices \(\operatorname{diag}(z,z)\) and
\(\operatorname{diag}(z^2,1)\). It gives the actual bundle isomorphism
\(L\oplus L\cong L^{\otimes2}\oplus\mathbf1\). Hence \(u^2=(1-[L])^2=0\)
integrally. The Thom isomorphism now gives \[
\begin{aligned}
K_c^0(L)&=\mathbb Z\xi\oplus\mathbb Z(\pi^*u\,\xi),\\
K_c^1(L)&=0.
\end{aligned}
\tag{12.3}
\] Zero-section restriction sends \(\xi\) to \(u\) and \(\pi^*u\,\xi\)
to \(u^2=0\). In a unitary fibre frame, (12.1) has comparison \(z\);
\href{supporting/relative-k-foundations/relative-k-foundations.html\#rk-planar-bott}{Theorem
RK.6 and Lemma RK.6a} prove its positive planar class and first Chern
value \(-1\) on that compactified fibre. The zero-section comparison
rules are
\href{supporting/relative-k-foundations/relative-k-foundations.html\#rk-symbol-rules}{Corollary
RK.4a}. Globally \[
\operatorname{ch}_c(\xi)
 =-\left(1-\frac12\pi^*h\right)U_L
\tag{12.4}
\] by Proposition 8.1. Thus the clutching sign, the zero-section
difference, and the Todd factor agree.

\subsection{What this lesson uses}\label{what-this-lesson-uses}

The Euclidean Fourier theorem and oscillator domain, compact graph
inclusion and kernel calculation are the exact earlier local results in
\emph{Bott periodicity in KK: the Bott and Dirac elements}, Lemmas 0.0a
and 4.0. That lesson proves its criterion locally and does not use the
present lesson, so this dependency has no cycle. The analytic
bounded-transform and product-recognition arguments are proved locally
in Lemmas 0.1--0.4 and Theorem 0.5, including local compactness and a
nonessential first representation. The central extension of the product
argument is proved in Lemma 1.1 above.

The bundle/projection and relative-triple correspondence, radial
identification, actual symbol rules and integral sphere groups are
proved in the earlier
\href{supporting/relative-k-foundations/relative-k-foundations.html\#rk-bundle-projection}{Relative
difference bundles, radial pairs and the planar normalization,
RK.1--RK.7}. The exact cohomological inputs used in Section 8 are proved
in
\href{supporting/cohomological-foundations/cohomological-foundations.html\#cf-thom}{Cohomological
foundations for the complex Thom comparison}: CF.0 supplies chains,
products and coefficients; CF.1b--CF.1c prove the ordinary Thom theorem
on every Hausdorff base and its orientation and product rules; CF.2
proves the quotient and radial comparisons; CF.3--CF.4 prove the
projective ring, first-Chern normalization, projective-bundle formula
and splitting injection; CF.5a--CF.5b prove the natural even relative
character, its bundle action and its actual symbol products; and CF.6
gives the complete outward-symbol Todd calculation. No full
odd-character or rational-isomorphism theorem is needed in Section 8,
and no singular-cohomology comparison isomorphism on all compact
Hausdorff spaces is asserted.

Lemma 9.0 proves the compact equivariant product and associativity
needed in Theorem 9.1 from the preceding non-equivariant connection,
partition and associativity proofs. The later lesson \emph{Equivariant
KK-theory and the Green--Julg theorem} develops stabilization,
crossed-product comparison and induction; none of those results is
needed here. The later lesson \emph{Wrong-way maps for K-oriented maps}
treats the further composition and independence assertions in the
embedding example.

The inward-to-outward vector comparison in \emph{A fundamental class for
an action on a manifold}, in \emph{Cyclic cohomology, connections and
transverse geometry}, Lemma 8.12, uses the same geometric Clifford
symbol. Its antipodal sign \((-1)^N\) is the parity sign of reversing
all fibre coordinates. It is separate from the product conversion
\((-1)^{rs}\) in Proposition 7.2. When an outward coordinate class is
normalized as an ordered product of positive line classes, translate
that statement into this course using Corollary 7.3: with
\(\boxtimes_R\) the class is \(\epsilon_N\tau\), while retaining
\(\tau\) uses \(\boxtimes_{\mathrm{out}}\). The planar projection in
that lesson is constant-unitarily conjugate to the planar projection
used here; that agreement of ordinary projection classes does not
identify the odd-odd product or its degree-two Morita convention. These
formulas provide the conversion needed before combining the two
constructions in an index argument.

\subsection{References}\label{references}

\begin{itemize}
\item
  B. Blackadar,
  \href{https://www.bruceblackadar.com/Mathematics/book6.pdf}{\emph{K-Theory
  for Operator Algebras}, author-posted corrected second edition}. The
  \href{https://www.bruceblackadar.com/mathpubs.html}{author identifies
  this freely downloadable version} as a slightly corrected second
  edition. The source retains its author copyright and the usage terms
  stated on that page. Sections 19.9.3--19.9.4 and 20.3.2, printed
  pp.~201--202 and 208, describe the Clifford, parametrized Thom and
  compact-equivariant constructions. Those passages include exercises
  and statements rather than complete proofs of all inverse identities.
  Sections 2--9 above prove the needed domain, product, sign, bundle and
  inverse assertions, using the exact earlier oscillator proof and the
  local analytic capsule in Section 0. Lemma 9.0 supplies the compact
  equivariant foundation before Theorem 9.1 uses it.
\item
  K. van den Dungen,
  \href{https://arxiv.org/pdf/2006.10616v2}{\emph{Localisations of
  half-closed modules and the unbounded Kasparov product}, arXiv version
  2}, 28 March 2022. The freely readable Sections 2.4 and 3.1, printed
  pp.~11--16, supply the source comparison for creation estimates and
  integrated positivity. Section 0 above proves the complete criterion
  needed in this lesson, including its nonessential representation case.
  No open adaptation licence is inferred from free access.
\end{itemize}

\end{document}
